\documentclass[aspectratio=169]{beamer}
\input{header}
\coursecode{JKSSB}

\title{Networking Basics}
\subtitle{Lesson 40 --- Networks, Topologies and Devices}
\author{JKSSB Computer Awareness}
\date{\today}

\begin{document}
\frame{\titlepage}

\begin{frame}{What Is a Computer Network?}
  \begin{itemize}
    \item A \key{computer network} is two or more computers (and other
      devices) connected so they can \key{share data and resources}.
    \item Shared resources include files, printers, an internet connection,
      and application software.
    \item Every participant in the network is called a \key{node}.
    \item A \key{host} is a node that can send or receive data over the
      network, such as a computer or a printer.
  \end{itemize}
  \begin{block}{The one idea}
    A network connects nodes so that data and resources can be shared.
  \end{block}
\end{frame}

\begin{frame}{Client and Server}
  \begin{itemize}
    \item A \key{client} is the device (or program) that \key{requests} a
      service, such as a web page or a file.
    \item A \key{server} is the computer that \key{provides} the service:
      it stores shared files or pages and answers client requests.
    \item In the common \key{client-server} arrangement, many clients ask and
      one server answers.
    \item \muted{The office desktop that opens a shared file is the client;
      the machine storing that file is the server.}
  \end{itemize}
  \begin{block}{Keep the pair straight}
    The client requests; the server provides. Do not swap the roles.
  \end{block}
\end{frame}

\begin{frame}{Network Types by Area: PAN, LAN, MAN, WAN}
  \begin{center}
    \begin{tabular}{p{0.10\textwidth}p{0.26\textwidth}p{0.50\textwidth}}
      \toprule
      \textbf{Type} & \textbf{Full form} & \textbf{Coverage} \\
      \midrule
      PAN & Personal Area Network & One person: phone, laptop, earphones \\
      LAN & Local Area Network & One building or campus: office, school \\
      MAN & Metropolitan Area Network & One city: branches linked across town \\
      WAN & Wide Area Network & Country or world: the Internet is the largest WAN \\
      \bottomrule
    \end{tabular}
  \end{center}
  \vspace{0.25cm}
  \muted{Area grows from PAN (smallest) up to WAN (largest).}
\end{frame}

\begin{frame}{WLAN, Intranet and the Internet}
  \begin{itemize}
    \item A \key{WLAN} (Wireless Local Area Network) is a LAN built with
      wireless signals instead of cables.
    \item An \key{intranet} is a private network inside one organisation;
      only its own staff can use it.
    \item The \key{Internet} is the worldwide public network of connected
      networks; anyone with access can use it.
    \item An \key{ISP} (Internet Service Provider) is the company that gives
      a customer access to the Internet.
  \end{itemize}
  \begin{alertblock}{Trap alert}
    An intranet is \key{private} to one organisation; the Internet is
    \key{public} and worldwide. Do not confuse the two words.
  \end{alertblock}
\end{frame}

\begin{frame}{Network Topologies: What the Word Means}
  \begin{itemize}
    \item A \key{topology} is the arrangement (layout) by which nodes are
      connected in a network.
    \item The five topologies in this lesson are \key{bus, star, ring, mesh,
      and tree}.
    \item Each topology decides how data travels and what happens when one
      cable or node fails.
    \item \muted{Exam questions ask: which topology has a central device?
      Which one fails fully when the main cable breaks?}
  \end{itemize}
  \begin{block}{How to remember them}
    Bus = one line; star = centre; ring = circle; mesh = every-to-every;
    tree = stars joined like branches.
  \end{block}
\end{frame}

\begin{frame}{Bus Topology}
  \begin{itemize}
    \item In a \key{bus} topology, all nodes share \key{one single main
      cable} (the bus).
    \item It is simple and cheap for a small network.
    \item Data sent by one node travels along the whole bus and is received
      by the addressed node.
    \item If the \key{main cable breaks}, the whole network fails.
  \end{itemize}
  \begin{alertblock}{Bus: one-line weakness}
    One shared cable is the strength (cheap) and the weakness (one break
    stops everything).
  \end{alertblock}
\end{frame}

\begin{frame}{Star Topology}
  \begin{itemize}
    \item In a \key{star} topology, every node connects to one
      \key{central device} (a hub or switch).
    \item All data passes through the centre.
    \item If one cable fails, only that node is affected; the rest keep
      working.
    \item If the \key{central device} fails, the whole network stops.
  \end{itemize}
  \begin{block}{Star is the office standard}
    The office LAN with one switch and many desktops is a star.
  \end{block}
\end{frame}

\begin{frame}{Ring Topology}
  \begin{itemize}
    \item In a \key{ring} topology, each node connects to exactly
      \key{two neighbours}, forming a closed circle.
    \item Data travels around the ring, usually in one direction, from node
      to node until it reaches its destination.
    \item Each node passes on (repeats) the data it is not addressed to.
    \item A break at one point disturbs the whole ring unless a second
      (dual) ring provides a backup path.
  \end{itemize}
\end{frame}

\begin{frame}{Mesh and Tree Topologies}
  \begin{columns}[T]
    \column{0.46\textwidth}
      \textbf{Mesh topology}
      \begin{itemize}
        \item Every node has a direct link to many (or all) other nodes.
        \item Very \key{reliable}: if one path fails, another path carries
          the data.
        \item Costly: it needs the most cable and ports.
      \end{itemize}
    \column{0.46\textwidth}
      \textbf{Tree topology}
      \begin{itemize}
        \item Several \key{star} networks joined through a central backbone,
          like branches of a tree.
        \item Suits a large campus: each department is a star, joined at
          the centre.
        \item If the backbone fails, whole branches lose connection.
      \end{itemize}
  \end{columns}
  \vspace{0.3cm}
  \muted{Mesh = most reliable and most costly; tree = stars joined together.}
\end{frame}

\begin{frame}{Topology Comparison at a Glance}
  \begin{center}
    \begin{tabular}{p{0.12\textwidth}p{0.30\textwidth}p{0.44\textwidth}}
      \toprule
      \textbf{Topology} & \textbf{Shape} & \textbf{Single point of failure} \\
      \midrule
      Bus & One shared cable & Main cable breaks: all fail \\
      Star & Centre + spokes & Central hub/switch fails: all fail \\
      Ring & Closed circle & One break disturbs the ring \\
      Mesh & Many direct links & No single break stops it \\
      Tree & Stars on a backbone & Backbone fails: branches fail \\
      \bottomrule
    \end{tabular}
  \end{center}
  \vspace{0.25cm}
  \muted{Read this table once and the matching items answer themselves.}
\end{frame}

\begin{frame}{Transmission Media: Guided and Unguided}
  \begin{itemize}
    \item \key{Transmission media} are the paths that carry data signals
      from one device to another.
    \item \key{Guided (wired) media} confine the signal inside a cable:
      twisted-pair cable, coaxial cable, and fibre-optic cable.
    \item \key{Unguided (wireless) media} carry the signal through air or
      space: radio waves, microwaves, satellite links, and infrared.
    \item Guided media need a physical cable; unguided media need antennas
      or line of sight.
  \end{itemize}
  \begin{block}{The one split}
    Guided = through a cable; unguided = through the air.
  \end{block}
\end{frame}

\begin{frame}{Guided Media: The Three Cables}
  \begin{itemize}
    \item \key{Twisted-pair cable}: pairs of copper wires twisted together;
      the common telephone and office-LAN cable. Cheap, short range.
    \item \key{Coaxial cable}: a central copper wire with insulation and a
      shield; used for cable TV and older LANs. Better shielding than
      twisted pair.
    \item \key{Fibre-optic cable}: thin glass or plastic strands carrying
      \key{light} signals; fastest, longest range, immune to electrical
      noise, but costliest.
  \end{itemize}
  \begin{alertblock}{Fibre is fastest}
    Rank by speed and range: fibre-optic is fastest, then coaxial, then
    twisted pair is slowest of the three.
  \end{alertblock}
\end{frame}

\begin{frame}{Unguided Media: Wireless Signals}
  \begin{itemize}
    \item \key{Radio waves}: travel in all directions; used for FM radio and
      WLAN/Bluetooth over short ranges.
    \item \key{Microwaves}: travel in a straight line (line of sight); used
      for point-to-point links between towers.
    \item \key{Satellite links}: microwave signals relayed through a
      satellite; cover very long distances but with delay.
    \item \key{Infrared}: very short range, blocked by walls; used for TV
      remotes.
  \end{itemize}
\end{frame}

\begin{frame}{Networking Devices I: NIC, Modem, Hub}
  \begin{itemize}
    \item \key{NIC} (Network Interface Card): the card (or built-in circuit)
      that lets a computer physically connect to a network.
    \item \key{Modem} (modulator-demodulator): converts digital computer
      signals to analogue signals for a telephone line and back again; it
      connects a home or office to the ISP.
    \item \key{Hub}: connects several nodes in a star and \key{broadcasts}
      every incoming signal to all ports; the addressed node accepts it.
  \end{itemize}
  \begin{block}{Hub vs everything later}
    A hub sends data to \key{everyone}; a switch sends it only to the
    addressed device.
  \end{block}
\end{frame}

\begin{frame}{Networking Devices II: Switch, Bridge, Repeater}
  \begin{itemize}
    \item \key{Switch}: connects nodes in a star and forwards each message
      \key{only to the addressed} port. Smarter and faster than a hub.
    \item \key{Bridge}: connects \key{two LAN segments} and filters traffic,
      passing only data meant for the other side.
    \item \key{Repeater}: \key{regenerates} a weak signal so it can travel
      further along a long cable; it works at the signal level only.
  \end{itemize}
  \begin{alertblock}{Trap alert}
    Hub = to all; switch = to the addressed port; repeater = only strengthens
    the signal, it does not choose a path.
  \end{alertblock}
\end{frame}

\begin{frame}{Networking Devices III: Router, Gateway, Access Point, Firewall}
  \begin{itemize}
    \item \key{Router}: connects \key{different networks} (for example a home
      LAN to the ISP) and \key{chooses the best path} for each packet.
    \item \key{Gateway}: connects networks that use \key{different protocols}
      or formats, translating between them.
    \item \key{Access point}: the wireless base station that lets Wi-Fi
      devices join a wired network.
    \item \key{Firewall}: \key{blocks unauthorised access} while letting
      legitimate traffic pass; a security device (hardware, software, or
      both).
  \end{itemize}
  \begin{block}{Router vs gateway vs firewall}
    Router chooses the path; gateway translates between different networks;
    firewall guards the entry.
  \end{block}
\end{frame}

\begin{frame}{High-Yield Drill: Device--Role Matching}
  \begin{center}
    \begin{tabular}{p{0.24\textwidth}p{0.62\textwidth}}
      \toprule
      \textbf{Device} & \textbf{One-line role} \\
      \midrule
      NIC & Connects a computer to the network \\
      Modem & Digital $\leftrightarrow$ analogue for the phone line / ISP link \\
      Hub & Broadcasts to all ports \\
      Switch & Forwards only to the addressed port \\
      Repeater & Regenerates a weak signal \\
      Bridge & Joins two LAN segments, filters traffic \\
      Router & Joins different networks, chooses the path \\
      Gateway & Translates between different protocols \\
      Access point & Lets wireless devices join the network \\
      Firewall & Blocks unauthorised access \\
      \bottomrule
    \end{tabular}
  \end{center}
  \vspace{0.25cm}
  \muted{Device--role matching is the highest-yield item type in this lesson.}
\end{frame}

\begin{frame}{Trap Alert: Networking Facts to Keep Straight}
  \begin{itemize}
    \item A hub sends data to \key{all} nodes; a \key{switch} sends it only
      to the addressed node.
    \item A \key{repeater} only regenerates the signal; the \key{router}
      chooses the path.
    \item A \key{bridge} joins two similar LAN segments; a \key{gateway}
      joins (translates between) different protocols.
    \item \key{LAN} is one building; \key{MAN} is one city; \key{WAN} spans
      countries; \key{PAN} is one person.
    \item \key{Intranet} is private; the \key{Internet} is public.
  \end{itemize}
\end{frame}

\begin{frame}{Lesson 40: Recall}
  \begin{itemize}
    \item A network links nodes to share data; a client requests and a
      server provides.
    \item By area: PAN (person), LAN (building), MAN (city), WAN (world);
      WLAN is wireless LAN; ISP = Internet Service Provider.
    \item Topologies: bus (one cable), star (centre), ring (circle), mesh
      (many links), tree (stars on a backbone).
    \item Guided media: twisted pair, coaxial, fibre-optic (fastest);
      unguided: radio, microwave, satellite, infrared.
    \item Devices: NIC connects; modem converts for the line; hub
      broadcasts; switch forwards to the address; repeater regenerates;
      bridge joins LAN segments; router paths between networks; gateway
      translates; access point admits Wi-Fi; firewall guards.
  \end{itemize}
\end{frame}
\end{document}
